% !TeX root = ../main.tex
% Lezione 5.pdf, pp. 2--11. La proprietà di Archimede precede gli esempi
% che la usano. La parte sulle funzioni, pp. 12--16, è nel capitolo 05.
% Corretto il passaggio del manoscritto su 1/n < -1+epsilon:
% per 0<epsilon<1 il secondo membro è negativo e non vi sono soluzioni.
\section{Il principio di Archimede nei reali}
% Fonte: Lezione 5, pp.~8--9

\begin{teorema}[Principio di Archimede]
\label{teo:archimede-reali}
Fissati $a,b\in\R$, con $a>0$, esiste $n\in\N$ tale che $na>b$.
\end{teorema}
\begin{proof}
\textbf{Riduzione al supremo dei multipli.} Fissiamo $a>0$ e consideriamo l'insieme
\[
  A=\{na:n\in\N\}.
\]
È non vuoto, perché $0\in A$. Dimostriamo che è superiormente illimitato. Se per assurdo fosse superiormente limitato, il teorema di esistenza degli estremi garantirebbe un numero reale $S=\sup A$.

\textbf{Il successivo supera il supremo.} La caratterizzazione operativa, con la scelta $\varepsilon=a>0$, darebbe un elemento $\bar n a\in A$ per cui
\[
  \bar n a>S-a.
\]
Sommando $a$ ai due membri otteniamo $(\bar n+1)a>S$. Ma $\bar n+1\in\N$, quindi anche $(\bar n+1)a\in A$, in contraddizione con il fatto che $S$ sia un maggiorante. L'insieme dei multipli è dunque superiormente illimitato: per ogni $b\in\R$ contiene un elemento $na>b$, come richiesto.
\end{proof}

\textbf{Conseguenze.}
\begin{itemize}
  \item Con $a=1$ si ottiene che i naturali superano qualunque reale: $\N$ è superiormente illimitato.
  \item I reciproci dei naturali positivi possono essere arbitrariamente piccoli:
  \[
    \forall\varepsilon>0\ \exists n\in\N^*:\quad \frac1n<\varepsilon.
  \]
  Basta infatti scegliere $n>1/\varepsilon$.
\end{itemize}
La completezza ha quindi una conseguenza concreta: non esiste un reale positivo minore o uguale a tutti i reciproci dei naturali positivi.

\addtocontents{toc}{\protect\Needspace{3\protect\baselineskip}}
\section{Esempi ed esercizi su supremo e infimo}
% Fonte: Lezione 5, pp.~2--7

\subsection{Un intervallo e un insieme con un punto isolato}
% Fonte: Lezione 5, p.~2

\begin{esempio}[Intervallo con estremo escluso]
Per $A=[0,1)$ si ha
\[
  \inf A=\min A=0,\qquad \sup A=1.
\]
\textbf{Minimo.} Il minimo è $0$, perché $0\in A$ e tutti gli elementi sono non negativi.

\textbf{Supremo non raggiunto.} Il numero $1$ è un maggiorante e, dato $\varepsilon>0$, il punto $a=1-\min\{\varepsilon/2,1/2\}$ appartiene ad $A$ e soddisfa $a>1-\varepsilon$. Il criterio operativo dimostra dunque che $1=\sup A$. Tuttavia $1\notin A$, quindi il massimo non esiste.
\end{esempio}

\begin{esempio}[Un massimo isolato]
Consideriamo
\[
  B=[-1,2]\cup\{5\}.
\]
In questo caso entrambi gli estremi sono raggiunti:
\[
  \inf B=\min B=-1,\qquad \sup B=\max B=5.
\]
Per verificare il criterio del supremo basta scegliere $a=5$ per ogni $\varepsilon>0$, poiché $5>5-\varepsilon$. Il punto che realizza il massimo può dunque essere isolato; la presenza dell'intero intervallo fra $2$ e $5$ non è richiesta.
\end{esempio}

\begin{schema}
\begin{tikzpicture}[x=1.45cm,y=1cm]
  \foreach \quota/\nome in {.7/A,-.8/B} {
    \draw[asse] (-1.5,\quota) -- (5.5,\quota);
    \node[anchor=east] at (-1.8,\quota) {$\nome$};
  }
  \draw[intervallo] (0,.7) -- (1,.7);
  \fill[coloreschema] (0,.7) circle[radius=2.5pt];
  \draw[estremo] (1,.7) circle[radius=2.5pt];
  \node[below=6pt] at (0,.7) {$0$};
  \node[below=6pt] at (1,.7) {$1$};
  \draw[intervallo] (-1,-.8) -- (2,-.8);
  \foreach \pos in {-1,2,5} {
    \fill[coloreschema] (\pos,-.8) circle[radius=2.5pt];
    \node[below=6pt] at (\pos,-.8) {$\pos$};
  }
\end{tikzpicture}
\par\smallskip In $A$ l'estremo superiore $1$ è escluso; in $B$ il punto $5$ è presente ed è il massimo, anche se è separato dall'intervallo.
\end{schema}

\subsection{I reciproci dei naturali positivi}
% Fonte: Lezione 5, pp.~3--4

\begin{esempio}[Reciproci dei naturali positivi]
Sia
\[
  A=\left\{\frac1n:n\in\N^*\right\}.
\]
La rappresentazione dei suoi punti suggerisce $\sup A=1$ e $\inf A=0$. Verifichiamo entrambe le affermazioni.

\textbf{Massimo e supremo.} Il numero $1$ appartiene ad $A$ per $n=1$ e maggiora tutti gli elementi, poiché $1/n\le1$ per ogni $n\ge1$. Dunque
\[
  1=\max A=\sup A.
\]
\textbf{Infimo non raggiunto.} Il numero $0$ non appartiene ad $A$: nessun reciproco di un naturale positivo è zero. Per dimostrare che è l'infimo usiamo le due condizioni operative. È un minorante, perché $1/n>0$; inoltre, per ogni $\varepsilon>0$, Archimede fornisce un naturale $n>1/\varepsilon$, e quindi
\[
  \frac1n<\varepsilon=0+\varepsilon.
\]
Ne segue $\inf A=0$. Poiché $0\notin A$, il minimo non esiste.
\end{esempio}

\subsection{Come scartare un candidato sbagliato per l'infimo}
% Fonte: Lezione 5, pp.~5--6

Per lo stesso insieme $A=\{1/n:n\in\N^*\}$ proviamo a proporre un candidato $k$ diverso da $0$.

\textbf{Candidato positivo: non è un minorante.} Se $k>0$, $k$ non è un minorante. Infatti esiste $n>1/k$, quindi $1/n<k$.

\textbf{Candidato negativo: manca l’avvicinabilità.} Se $k<0$, esso è un minorante, perché tutti gli elementi di $A$ sono positivi, ma non è il più grande: $0$ è un minorante maggiore di $k$. Lo stesso difetto si vede con il criterio operativo. Per il candidato $k=-1$, prendendo $0<\varepsilon<1$ avremmo bisogno di un elemento per cui
\[
  \frac1n<-1+\varepsilon<0,
\]
che è impossibile. In generale, per un qualunque $k<0$, basta scegliere $\varepsilon=-k/2$: allora $k+\varepsilon=k/2<0$ e nessun elemento di $A$ può trovarsi sotto questa soglia.

\subsection{Due esercizi: cubi traslati e unione di insiemi}
% Fonte: Lezione 5, p.~7

\begin{esercizio}[Cubi traslati]
Determinare supremo e infimo dell'insieme
\[
  A_1=\{n^3-1:n\in\N\}.
\]
Stabilire se esistono massimo e minimo.
\end{esercizio}
\begin{svolgimento}
Poiché $n^3\ge0$ per ogni $n\in\N$ e il valore $-1$ si ottiene per $n=0$, si ha $\inf A_1=\min A_1=-1$.

\textbf{Illimitatezza superiore.} L'insieme è superiormente illimitato: dato un qualunque $M\in\R$, scegliamo un naturale $n>\max\{1,M+1\}$. Allora $n^3-1\ge n-1>M$. Dunque, nella retta estesa,
\[
  \sup A_1=+\infty.
\]
L'illimitatezza superiore esclude l'esistenza di un massimo.
\end{svolgimento}

\begin{esercizio}[Unione di due insiemi]
Determinare supremo, infimo ed eventuali massimo e minimo dell'insieme $A_2=\{x\in\R:x^2<1\}\cup\{x\in\R:(x-3)^2=0\}$.
\end{esercizio}
\begin{svolgimento}
\textbf{Descrizione dell'insieme.} Si ha
\[
  A_2=\{x\in\R:x^2<1\}
      \cup\{x\in\R:(x-3)^2=0\}
     =(-1,1)\cup\{3\}.
\]
Il numero $3$ appartiene ad $A_2$ e maggiora tutti gli elementi, quindi $\sup A_2=\max A_2=3$.

\textbf{Infimo non raggiunto.} Il numero $-1$ è un minorante. Dato $\varepsilon>0$, scegliamo
\[
  a=-1+\frac{\varepsilon}{1+\varepsilon}.
\]
Il punto $a$ appartiene a $(-1,0)\subseteq A_2$ e soddisfa $a<-1+\varepsilon$, perché $0<\varepsilon/(1+\varepsilon)<\varepsilon$. Perciò $\inf A_2=-1$; il minimo non esiste, perché $-1\notin A_2$.
\end{svolgimento}

% Mantiene insieme la dimostrazione e il disegno della densità.
\Needspace{36\baselineskip}
\section{La densità dei razionali nella retta reale}
% Fonte: Lezione 5, pp.~10--11

Fra due razionali distinti si trovano altri razionali. Il risultato seguente estende la proprietà a qualunque intervallo aperto con estremi reali.

\begin{teorema}[Densità di $\Q$ in $\R$]
\label{teo:densita-razionali}
Se $a,b\in\R$ e $a<b$, esiste $r\in\Q$ tale che $a<r<b$.
\end{teorema}
\begin{proof}
\textbf{Dilatazione dell’intervallo.} Poiché $b-a>0$, il principio di Archimede permette di scegliere $q\in\N^*$ tale che
\[
  q(b-a)>1.
\]
L'intervallo $(qa,qb)$ ha dunque ampiezza maggiore di $1$ e contiene un intero $p$.

\textbf{Scelta dell’intero.} Per ogni reale $t$ esiste un intero $p$ con $t<p\le t+1$. Infatti, scegliamo un naturale $N>-t$ e consideriamo il minimo naturale $m>t+N$, che esiste per Archimede e per il principio del minimo intero. Si ha $m\ge1$ e $m-1\le t+N$, altrimenti $m$ non sarebbe il minimo. Ponendo $p=m-N$ si ottiene proprio $t<p\le t+1$.

\textbf{Ritorno alla scala iniziale.} Applicando il fatto appena verificato a $t=qa$, troviamo
\[
  qa<p\le qa+1<qb.
\]
Dividendo per $q>0$ segue $a<p/q<b$. Il numero $r=p/q$ è razionale, perché $p\in\Z$ e $q\in\N^*$, e appartiene all'intervallo richiesto.
\end{proof}

\begin{schema}
\begin{tikzpicture}[x=1.6cm,y=1cm]
  \draw[asse] (-.3,0) -- (4.4,0);
  \draw[intervallo] (1.35,0) -- (3.25,0);
  \draw[estremo] (1.35,0) circle[radius=2.5pt];
  \draw[estremo] (3.25,0) circle[radius=2.5pt];
  \node[above=8pt,coloreschema] at (1.35,0) {$qa$};
  \node[above=8pt,coloreschema] at (3.25,0) {$qb$};
  \foreach \pos/\testo in {0/{p-2},1/{p-1},2/p,3/{p+1},4/{p+2}} {
    \draw[black!55] (\pos,-.08) -- (\pos,.08);
    \node[below=7pt] at (\pos,0) {$\testo$};
  }
  \fill[coloreschema] (2,0) circle[radius=2.5pt];
\end{tikzpicture}
\par\smallskip Dilatando l'intervallo con un fattore intero $q$ si ottiene un'ampiezza maggiore di $1$, che assicura la presenza di un intero $p$. Tornando alla scala iniziale si trova il razionale $p/q$.
\end{schema}

\textbf{Infinità dei punti razionali.} I razionali in $(a,b)$ sono infiniti: trovato un primo razionale $r_1$ nell'intervallo, si può cercarne un secondo in $(a,r_1)$, un terzo fra $a$ e il secondo, e continuare. Ogni intervallo aperto non vuoto della retta incontra quindi $\Q$, anche quando i suoi estremi sono irrazionali.
